\section*{Hoofdstuk \ref{ch:b1:det} --- Determinanten en lineaire
stelsels}

\begin{solution}{exo:b1:det:1}
$3 \times 2 - 1 \times 5 = 1$.

Tweede: $L_2 \leftarrow L_2 - L_1$, $L_3 \leftarrow L_3 - L_2$ (op
de oorspronkelijke rijen) geven de rijen $(1,2,3), (3,3,3),
(3,3,3)$: twee gelijke rijen, determinant $0$. (Sarrus bevestigt:
$45 + 84 + 96 - 105 - 48 - 72 = 0$.)

Derde: het is Vandermonde met $x = 1, 2, 3$
(\cref{ex:b1:det:vandermonde}): $(2-1)(3-1)(3-2) = 2$.
\end{solution}

\begin{solution}{exo:b1:det:2}
De determinant is gelijk aan (tel alle kolommen bij de eerste op en
splits af) $(\lambda + 2)$ maal
\[
\begin{vmatrix}
1 & 1 & \lambda\\ 1 & \lambda & 1\\ 1 & 1 & 1
\end{vmatrix}
= -(\lambda - 1)^2
\]
(veeg schoon met $L_1 \leftarrow L_1 - L_3$, $L_2 \leftarrow L_2 -
L_3$ en ontwikkel), wat $\det = -(\lambda+2)(\lambda-1)^2$ geeft.
\omterm{def:b1:vspaces:free}{Basis} $\iff \det \neq 0 \iff \lambda \notin \{1, -2\}$.
\end{solution}

\begin{solution}{exo:b1:det:3}
Eerste stelsel: $\det = -7$;
$x = \frac{1}{-7}\begin{vmatrix} 5 & 1\\ 4 & -2\end{vmatrix}
= \frac{-14}{-7} = 2$,
$y = \frac{1}{-7}\begin{vmatrix} 2 & 5\\ 3 & 4\end{vmatrix}
= \frac{-7}{-7} = 1$. Controle: $2(2) + 1 = 5$; $3(2) - 2 = 4$.

Tweede stelsel: na $L_2 - L_1$ en $L_3 - 2L_1$ worden de rijen
$(1,1,1)$, $(0,-2,0)$, $(0,-1,-3)$, dus
\[
\det A = \begin{vmatrix} 1&1&1\\ 1&-1&1\\ 2&1&-1\end{vmatrix}
= 1 \times \begin{vmatrix} -2 & 0\\ -1 & -3\end{vmatrix} = 6 .
\]
Cramer, met de kolommen vervangen door $(6,2,1)^{\mathsf T}$:
\[
x = \frac{6}{6} = 1, \qquad
y = \frac{12}{6} = 2, \qquad
z = \frac{18}{6} = 3
\]
(de tellers op dezelfde wijze berekend). Controle: $1 + 2 + 3 = 6$;
$1 - 2 + 3 = 2$; $2 + 2 - 3 = 1$.
\end{solution}

\begin{solution}{exo:b1:det:4}
Reduceer de aangevulde matrix: $L_2 \leftarrow L_2 - 2L_1$, $L_3
\leftarrow L_3 - L_1$:
\[
\begin{pmatrix}
1 & 2 & -1 & 1 & 1\\
0 & 0 & 3 & -3 & 3\\
0 & 0 & 3 & -3 & 3
\end{pmatrix}
\to
\begin{pmatrix}
1 & 2 & -1 & 1 & 1\\
0 & 0 & 1 & -1 & 1\\
0 & 0 & 0 & 0 & 0
\end{pmatrix}.
\]
Spilonbekenden $x, z$; \omterm{def:b1:vspaces:free}{vrije} onbekenden $y, t$. Terugsubstitutie:
$z = 1 + t$, $x = 1 - 2y + z - t = 2 - 2y$.
Oplossingsverzameling:
\[
\{(2 - 2y,\; y,\; 1 + t,\; t) : y, t \in \R\}
= (2, 0, 1, 0) + \operatorname{Vect}\bigl((-2,1,0,0),\,
(0,0,1,1)\bigr),
\]
een affien vlak (dimensie $2 = 4 - \operatorname{rk} 2$) van $\R^4$.
\end{solution}

\begin{solution}{exo:b1:det:5}
Inductie; $n = 1$ is het lege product $= 1$. Voer voor de stap $C_k
\leftarrow C_k - x_1 C_{k-1}$ uit voor $k = n, n-1, \dots, 2$ (in
die volgorde, zodat elke bewerking een nog niet gewijzigde kolom
gebruikt). De eerste rij wordt $(1, 0, \dots, 0)$; in rij $i \geq 2$
wordt de $k$-de ingang $x_i^{k-1} - x_1 x_i^{k-2} = x_i^{k-2}(x_i -
x_1)$. Ontwikkelen langs de eerste rij en $(x_i - x_1)$ uit elke rij
$i$ afsplitsen geeft
\[
V(x_1, \dots, x_n)
= \prod_{i=2}^{n} (x_i - x_1)\cdot V(x_2, \dots, x_n),
\]
en de inductiehypothese vervolledigt het product $\prod_{i<j}(x_j -
x_i)$.
\end{solution}

\begin{solution}{exo:b1:det:6}
Ontwikkeling van $D_n$ langs de eerste rij: $D_n = 2 D_{n-1} -
1\cdot\begin{vmatrix} 1 & \ast\\ 0 & D_{n-2}\text{-blok}
\end{vmatrix}$; de tweede determinant is, ontwikkeld langs zijn
eerste kolom, gelijk aan $D_{n-2}$. Dus $D_n = 2D_{n-1} - D_{n-2}$,
dat wil zeggen $D_n - D_{n-1} = D_{n-1} - D_{n-2}$: de verschillen
zijn constant, gelijk aan $D_2 - D_1 = 1$. Bijgevolg is $D_n = D_1 +
(n - 1) = n + 1$. (Controle: $D_2 = 3$, en het geval $3\times3$ is
\cref{ex:b1:det:cofactor}: $D_3 = 4$.)
\end{solution}

\begin{solution}{exo:b1:det:7}
$C_1 \leftarrow C_1 + C_2 + C_3$ maakt de eerste kolom constant
gelijk aan $(m+2)$; splits die af:
\[
\det = (m+2)\begin{vmatrix}
1 & 1 & m\\ 1 & m & 1\\ 1 & 1 & 1
\end{vmatrix}
\overset{L_1 - L_3,\ L_2 - L_3}{=}
(m+2)\begin{vmatrix}
0 & 0 & m-1\\ 0 & m-1 & 0\\ 1 & 1 & 1
\end{vmatrix}
= (m+2)\cdot\bigl(-(m-1)^2\bigr)
\]
(ontwikkel langs de eerste kolom: de enige ingang $1$ draagt teken
$+$, en de resterende $2 \times 2$ determinant is $0 \cdot 0 -
(m-1)(m-1) = -(m-1)^2$).

Bespreking. $m \notin \{1, -2\}$: een stelsel van Cramer; wegens de
symmetrie van de vergelijkingen is $x = y = z$, en geeft elke
vergelijking $(m + 2)x = 1$: eenduidige oplossing
$\bigl(\frac{1}{m+2}, \frac{1}{m+2}, \frac{1}{m+2}\bigr)$. $m = 1$:
alle drie de vergelijkingen luiden $x + y + z = 1$: de oplossingen
vormen het affiene vlak $x + y + z = 1$. $m = -2$: het optellen van
de drie vergelijkingen geeft $0 = 3$: lege
oplossingsverzameling.
\end{solution}

\begin{solution}{exo:b1:det:8}
($\Rightarrow$) Heeft $A^{-1}$ gehele ingangen, dan is $\det A \cdot
\det A^{-1} = 1$ met beide determinanten geheel (sommen van
producten van ingangen): twee gehele getallen met product $1$ zijn
beide $\pm1$.

($\Leftarrow$) De cofactorformule $A^{-1} = \frac{1}{\det A}
\operatorname{Com}(A)^{\mathsf T}$ (voor $n \leq 3$ door
rechtstreeks ontwikkelen nagegaan, in het algemeen aangenomen) heeft
een $\operatorname{Com}(A)$ met gehele ingangen (elke cofactor is
een gehele determinant); delen door $\det A = \pm 1$ behoudt de
geheeltalligheid.
\end{solution}

\begin{solution}{exo:b1:det:9}
Tel alle kolommen bij de eerste op: elke ingang van de nieuwe eerste
kolom is $a + nb$; splits die af, zodat de eerste kolom uit enen
bestaat. De \omterm{met:b1:matrices:gauss}{rij-operaties} $L_i \leftarrow L_i - L_1$ ($i \geq 2$)
vegen dan elke ingang onder de $1$ linksboven schoon en laten $a$ op
de diagonaal en $0$ elders in die rijen: de matrix is een bovenste
driehoeksmatrix met diagonaal $(1, a, \dots, a)$. Bijgevolg
\[
\det(aI + bJ) = (a + nb)\, a^{\,n-1} ,
\]
ongelijk aan nul dan en slechts dan als $a \neq 0$ en $a + nb \neq
0$: de voorwaarde van \cref{exo:b1:matrices:9}.
\end{solution}

\begin{solution}{exo:b1:det:10}
Blokoperaties (elk een samenstelling van de overeenkomstige $n$
scalaire bewerkingen, toegestaan door
\cref{thm:b1:det:props}~(1)):
\[
\begin{vmatrix} A & B\\ B & A\end{vmatrix}
\overset{C_1 \leftarrow C_1 + C_2}{=}
\begin{vmatrix} A + B & B\\ A + B & A\end{vmatrix}
\overset{L_2 \leftarrow L_2 - L_1}{=}
\begin{vmatrix} A + B & B\\ 0 & A - B\end{vmatrix}
= \det(A+B)\,\det(A-B),
\]
met de blokdriehoeksregel voor de laatste stap.
\end{solution}

\begin{solution}{exo:b1:det:11}
$C_1 \leftarrow C_1 + C_2 + C_3$ maakt de eerste kolom constant
gelijk aan $(a + b + c)$; splits die af, en daarna $L_2 \leftarrow
L_2 - L_1$, $L_3 \leftarrow L_3 - L_1$:
\[
\Delta = (a+b+c)\begin{vmatrix}
1 & b & c\\
0 & a - b & b - c\\
0 & c - b & a - c
\end{vmatrix}
= (a+b+c)\bigl[(a-b)(a-c) + (b-c)^2\bigr],
\]
en uitgewerkt is $(a-b)(a-c) + (b-c)^2 = a^2 + b^2 + c^2 - ab - bc
- ca$. Over $\C$, met $j^3 = 1$ en $1 + j + j^2 = 0$:
\begin{align*}
(a + jb + j^2c)(a + j^2b + jc)
&= a^2 + b^2 + c^2 + (j + j^2)(ab + bc + ca)\\
&= a^2 + b^2 + c^2 - ab - bc - ca ,
\end{align*}
waaruit de volledige ontbinding volgt. (Structureel: de kolom $(1,
j, j^2)^{\mathsf T}$ voldoet aan $M\,(1, j, j^2)^{\mathsf T} = (a +
jb + j^2c)(1, j, j^2)^{\mathsf T}$, en net zo voor $j^2$ en $1$: de
drie factoren zijn de drie ``eigenwaarden'' van de circulant, een
verhaal dat in het volume van bachelorjaar 2 wordt
gesystematiseerd.)
\end{solution}

\begin{solution}{exo:b1:det:12}
Schrijf $r = \operatorname{rk} A$.

\emph{Er bestaat een inverteerbare $r \times r$ deelmatrix.} Kies
$r$ \omterm{def:b1:vspaces:free}{vrije} kolommen van $A$ en zij $B \in \mathcal{M}_{n,r}$ de
matrix die zij vormen: $\operatorname{rk} B = r$. Omdat de rijrang
gelijk is aan de kolomrang (\cref{thm:b1:matrices:rank}), heeft $B$
$r$ \omterm{def:b1:vspaces:free}{vrije} rijen; die rijen behouden levert een $r \times r$
deelmatrix van $A$ van rang $r$, dus inverteerbaar.

\emph{Een grotere bestaat niet.} Zij $S$ een inverteerbare $s \times
s$ deelmatrix, genomen uit de kolommen $j_1, \dots, j_s$ en de
rijen $i_1, \dots, i_s$ van $A$. Wordt een combinatie $\sum_k
\lambda_k C_{j_k} = 0$ van de bijbehorende \emph{volledige} kolommen
nul, dan geeft alleen de rijen $i_1, \dots, i_s$ lezen dat $\sum_k
\lambda_k S_k = 0$ op de kolommen van $S$, dus zijn alle $\lambda_k
= 0$ ($S$ is inverteerbaar): de kolommen $C_{j_1}, \dots, C_{j_s}$
van $A$ zijn \omterm{def:b1:vspaces:free}{vrij}, en $s \leq \operatorname{rk} A = r$.

Bijgevolg is $\operatorname{rk} A$ precies de grootste omvang van een
inverteerbare deelmatrix.
\end{solution}

\begin{solution}{pb:b1:det:1}
\textbf{1.} $V(1,2,3,4) = (2-1)(3-1)(4-1)(3-2)(4-2)(4-3) = 1
\cdot 2\cdot 3\cdot 1\cdot 2\cdot 1 = 12$. Interpolatie in
verschillende knooppunten vraagt om de coëfficiënten van $P$ die
$W c = y$ oplossen met $W = (x_i^{\,j-1})$, en $\det W = V \neq 0$:
een stelsel van Cramer.

\textbf{2.} Werk de kolommen van links naar rechts af. $C_1$ is de
constante kolom $P_0(x_i) = 1$ ($P_0$ is \omterm{def:b1:poly:def}{monisch} van graad $0$).
Neem aan dat de kolommen $1, \dots, j-1$ al zijn gereduceerd tot de
zuivere machten $1, x_i, \dots, x_i^{\,j-2}$. Omdat $P_{j-1} =
X^{j-1} + \sum_{k < j-1}\alpha_k X^k$, laat het aftrekken van de
combinatie $\sum_k \alpha_k\,(\text{kolom van } x_i^k)$ van $C_j$
--- een bewerking die de determinant niet verandert --- de zuivere
machtskolom $x_i^{\,j-1}$ over. Na de laatste kolom is de matrix de
Vandermonde-matrix: $\det = V(x_1, \dots, x_n)$.

\textbf{3.} De \omterm{def:b1:poly:def}{veeltermen} $(j-1)!\,B_{j-1}$ zijn \omterm{def:b1:poly:def}{monisch} van graad
$j - 1$, dus geeft vraag 2
\[
\det\bigl(B_{j-1}(m_i)\bigr)
= \frac{V(m_1, \dots, m_n)}{0!\,1!\cdots(n-1)!} .
\]
Het linkerlid is de determinant van een matrix met \emph{gehele}
ingangen ($B_k$ is geheelwaardig op $\Z$: de vragen 16--17 van de
weekendopgave \cref{pb:b1:vspaces:1}), dus een geheel getal; en het
is positief omdat $V(m_1, \dots, m_n) > 0$ voor $m_1 < \dots <
m_n$. De superfaculteit \omterm{def:b1:arith:divides}{deelt} dus het product van alle paarsgewijze
verschillen.

\textbf{4.} Splits $x_i$ af uit elke rij $i$:
$\det(x_i^{\,j})_{j = 1..n} = x_1\cdots x_n\,
\det(x_i^{\,j-1}) = x_1\cdots x_n\,V$.

\textbf{5.} $(W^{\mathsf T}W)_{ij} = \sum_k x_k^{\,i-1}
x_k^{\,j-1} = p_{i+j-2}$: $S = W^{\mathsf T}W$. Bijgevolg is $\det
S = \det(W^{\mathsf T})\det W = V^2$
(\cref{thm:b1:det:props}~(2),(4)). Voor reële $x_i$ is $\det S =
V^2 \geq 0$, en is $S$ inverteerbaar dan en slechts dan als $V \neq
0$, dus dan en slechts dan als de $x_i$ paarsgewijs verschillend
zijn --- een tekenbepaalde toets, berekenbaar uit de machtsommen
alleen.

\textbf{6.} De interpolatievoorwaarden $\sum_{k} c_k\,x_i^{\,k} =
y_i$ vormen het stelsel $Wc = y$; $\det W = V \neq 0$ geeft het
bestaan en de eenduidigheid in één keer. Dit is het derde bewijs in
het boek: de expliciete formule in \cref{thm:b1:poly:lagrange}, het
argument met de kern in \cref{ex:b1:linmaps:interpolation}, en
Cramer hier.

\textbf{7.} Cramer: $c_{n-1} = \det W'/\det W$ met $W'$ gelijk aan
$W$ met haar laatste kolom vervangen door $y$. Ontwikkeling van
$\det W'$ langs die kolom geeft
\[
\det W' = \sum_{i=1}^n (-1)^{i+n} y_i\,V(x_1, \dots, \widehat{x_i},
\dots, x_n) .
\]
Nu is $V = V(\setminus i)\cdot\prod_{j<i}(x_i - x_j)\prod_{j>i}
(x_j - x_i)$, en het omzetten van het tweede product kost
$(-1)^{n-i}$:
\[
(-1)^{i+n}\,\frac{V(\setminus i)}{V}
= \frac{(-1)^{i+n}(-1)^{n-i}}{\prod_{j\neq i}(x_i - x_j)}
= \frac{1}{\prod_{j\neq i}(x_i - x_j)} ,
\]
waaruit $c_{n-1} = \sum_i y_i/\prod_{j \neq i}(x_i - x_j)$ ---
opnieuw de formule van de \omterm{pb:b1:vspaces:1}{gedeelde differenties}.

\textbf{8.} $L_3 \leftarrow L_3 - L_1$ geeft de rijen $(1, x_1,
x_1^2)$, $(0, 1, 2x_1)$, $(0,\ x_2 - x_1,\
(x_2-x_1)(x_2+x_1))$; ontwikkeling langs de eerste kolom en het
afsplitsen van $(x_2 - x_1)$ geven
\[
(x_2 - x_1)\begin{vmatrix} 1 & 2x_1\\ 1 & x_2 + x_1
\end{vmatrix} = (x_2 - x_1)(x_2 - x_1) = (x_2 - x_1)^2 .
\]
Ongelijk aan nul voor $x_1 \neq x_2$: het \omterm{def:b1:linmaps:def}{lineaire} stelsel dat
$P(x_1) = u$, $P'(x_1) = v$, $P(x_2) = w$ uitdrukt in de
coëfficiënten van $P \in \R_2[X]$ is een stelsel van Cramer --- de
interpolatie van Hermite met een verdubbeld knooppunt is goed
gesteld.

\textbf{9.} $P = a + bX + cX^2$ met $a = P(0) = 1$, $b = P'(0) =
0$, $a + b + c = P(1) = 2$: $c = 1$, dus $P = 1 + X^2$, eenduidig.
Consistentie: hier is $x_1 = 0$, $x_2 = 1$ en is de determinant van
vraag 8 gelijk aan $(1 - 0)^2 = 1 \neq 0$.

\textbf{10.} Rechtstreekse berekening:
\[
C_2 = \frac{1}{(a_1+b_1)(a_2+b_2)} - \frac{1}{(a_1+b_2)(a_2+b_1)}
= \frac{(a_1+b_2)(a_2+b_1) - (a_1+b_1)(a_2+b_2)}
{\prod_{i,j}(a_i+b_j)} ,
\]
en de teller werkt uit tot $a_1b_1 + a_2b_2 - a_1b_2 - a_2b_1 =
(a_2 - a_1)(b_2 - b_1)$: verschillen gedeeld door sommen.

\textbf{11.} Voor $i < n$ is de nieuwe ingang van rij $i$
\[
\frac{1}{a_i + b_j} - \frac{1}{a_n + b_j}
= \frac{a_n - a_i}{(a_i + b_j)(a_n + b_j)} .
\]
Splits $(a_n - a_i)$ af uit elke rij $i < n$, en daarna
$\frac1{a_n + b_j}$ uit elke kolom $j$: wat overblijft heeft
ingangen $\frac1{a_i + b_j}$ in de rijen $i < n$ en constant $1$ in
rij $n$ --- de matrix $M$, met de aangekondigde voorfactor.

\textbf{12.} Op $M$ maakt de bewerking $C_j \leftarrow C_j - C_n$
voor $j < n$ van rij $n$ de rij $(0, \dots, 0, 1)$, en in rij $i <
n$
\[
\frac{1}{a_i + b_j} - \frac{1}{a_i + b_n}
= \frac{b_n - b_j}{(a_i + b_j)(a_i + b_n)} .
\]
Splits $(b_n - b_j)$ af uit elke kolom $j < n$ en $\frac1{a_i +
b_n}$ uit elke rij $i < n$, en ontwikkel daarna langs de laatste rij
(teken $(-1)^{n+n} = +1$): de overblijvende determinant is
$C_{n-1}$. Het verzamelen van de factoren uit de vragen 11--12
geeft
\[
C_n = \frac{\prod_{i<n}(a_n - a_i)\,\prod_{j<n}(b_n - b_j)}
{\prod_{j}(a_n + b_j)\,\prod_{i<n}(a_i + b_n)}\;C_{n-1},
\]
en de inductie (\omterm{def:b1:vspaces:free}{basis} $C_1 = \frac1{a_1+b_1}$) zet precies de
dubbele alternant van Cauchy in elkaar: de factoren $(a_j -
a_i)(b_j - b_i)$ voor alle paren, gedeeld door alle sommen $(a_i +
b_j)$.

\textbf{13.} De formule wordt nul dan en slechts dan als een $a_j =
a_i$ of een $b_j = b_i$: de Cauchy-matrix is inverteerbaar dan en
slechts dan als beide families paarsgewijs verschillend zijn.
Hilbert: $a_i = i$, $b_j = j - 1$. Voor $n = 2$: teller $(2-1)(1-0)
= 1$, noemer $1\cdot2\cdot2\cdot3 = 12$: $\det H_2 = \frac1{12}$.
Voor $n = 3$: teller $\bigl[(1)(2)(1)\bigr]^2 = 4$, noemer
$(1\cdot2\cdot3)(2\cdot3\cdot4)(3\cdot4\cdot5) = 6\cdot24\cdot60 =
8640$: $\det H_3 = \frac{4}{8640} = \frac1{2160}$. Inverse voor $n =
2$:
\[
H_2^{-1} = 12\begin{pmatrix} \frac13 & -\frac12\\[2pt]
-\frac12 & 1\end{pmatrix}
= \begin{pmatrix} 4 & -6\\ -6 & 12 \end{pmatrix},
\]
allemaal gehele getallen (een verschijnsel dat voor elke $n$ waar
is).

\textbf{14.} De matrix van het stelsel is de Cauchy-matrix, volgens
vraag 13 inverteerbaar wanneer de $b_j$ (en de $a_i$) paarsgewijs
verschillend zijn: een eenduidige oplossing. Interpretatie: een
rationale functie $R = \sum_j \frac{c_j}{X + b_j}$ met
enkelvoudige \omterm{def:b1:fractions:field}{polen} wordt bepaald door $n$ van haar waarden
$R(a_1), \dots, R(a_n)$, en omgekeerd wordt elk zo'n gegevensblad
precies één keer gerealiseerd --- de bemonsteringstegenhanger van
de stelling over het bestaan en de eenduidigheid van de splitsing
in partieelbreuken (\cref{thm:b1:fractions:complex}).

\textbf{15.} Vieta voor $X^3 + pX + q$: $\lambda_1 + \lambda_2 +
\lambda_3 = 0$, $\sum_{i<j}\lambda_i\lambda_j = p$, dus $p_1 = 0$
en $p_2 = p_1^2 - 2p = -2p$. Elke wortel voldoet aan $\lambda^3 =
-p\lambda - q$; sommeren geeft $p_3 = -p\,p_1 - 3q = -3q$.
Vermenigvuldigen met $\lambda$ en sommeren geeft $p_4 = -p\,p_2 -
q\,p_1 = 2p^2$.

\textbf{16.} Volgens vraag 5 (de identiteit $S = W^{\mathsf T}W$ en
$\det S = V^2$ gelden over $\C$) is
\[
V^2 = \begin{vmatrix}
3 & 0 & -2p\\
0 & -2p & -3q\\
-2p & -3q & 2p^2
\end{vmatrix}
= 3\bigl(-4p^3 - 9q^2\bigr) + (-2p)\bigl(0 - 4p^2\bigr)
= -4p^3 - 27q^2 ,
\]
ontwikkeld langs de eerste rij.

\textbf{17.} Een herhaalde wortel betekent twee gelijke
$\lambda_i$, dus $V = 0$, dus $\operatorname{disc} = -4p^3 - 27q^2
= 0$. Voor $X^3 - 3X + 2$: $4(-3)^3 + 27\cdot4 = -108 + 108 = 0$,
in overeenstemming met de dubbele wortel $1$ van $(X-1)^2(X+2)$.

\textbf{18.} De niet-reële wortels van een reële
derdegraadsveelterm komen in toegevoegde paren, dus doen zich bij
$\operatorname{disc} \neq 0$ precies twee gevallen voor. Drie
verschillende reële wortels: $V$ is reëel en ongelijk aan nul, dus
$\operatorname{disc} = V^2 > 0$. Eén reële wortel $\lambda_1$ en
$\lambda_3 = \conj{\lambda_2} \notin \R$: dan is
\[
(\lambda_2 - \lambda_1)(\lambda_3 - \lambda_1) =
\abs{\lambda_2 - \lambda_1}^2 > 0,
\qquad
\lambda_3 - \lambda_2 = -2\iu\,\operatorname{Im}\lambda_2 \neq 0,
\]
zodat $V$ een zuiver imaginair getal ongelijk aan nul is en
$\operatorname{disc} = V^2 < 0$. De twee tekens karakteriseren de
twee gevallen.

\textbf{19.} Neem $b_i = a_i$ in de dubbele alternant: de teller is
$\prod_{i<j}(a_j - a_i)^2 > 0$ en de noemer $\prod_{i,j}(a_i + a_j)
> 0$ (alle ingangen positief): de determinant is positief. (In
latere taal: de kern $\frac1{x+y}$ is positief definiet.)

\textbf{20.} Volgens de vragen 2--3 is de determinant gelijk aan
$V(2,4,7)/(0!\,1!\,2!) = \frac{(4-2)(7-2)(7-4)}{2} = \frac{30}{2} =
15$. Rechtstreeks is de matrix
\[
\begin{pmatrix}
1 & 2 & 1\\
1 & 4 & 6\\
1 & 7 & 21
\end{pmatrix},
\qquad
\det = (84 - 42) - 2(21 - 6) + (7 - 4) = 42 - 30 + 3 = 15 .
\]

\textbf{21.} Stel $\sum_i c_i\,(\lambda_i^{\,k})_{k} = 0$ als rij.
Aflezen bij $k = 0, 1, \dots, n-1$ geeft $W^{\mathsf T}c = 0$ met
$W = (\lambda_i^{\,j-1})$ inverteerbaar ($\det = V \neq 0$, want de
$\lambda_i$ zijn verschillend): $c = 0$. De meetkundige rijen zijn
\omterm{def:b1:vspaces:free}{vrij}.

\textbf{22.} $a = b = (1, 2, 3)$: teller $\bigl[(2-1)(3-1)
(3-2)\bigr]^2 = 4$; noemer $\prod_{i,j}(i + j) =
(2\cdot3\cdot4)(3\cdot4\cdot5)(4\cdot5\cdot6) = 24\cdot60\cdot120 =
172800$. Bijgevolg is $\det\bigl(\frac1{i+j}\bigr) =
\frac{4}{172800} = \frac1{43200}$.

\textbf{23.} Is $x_i = x_j$, dan houdt de verwisseling van de twee
veranderlijken het punt vast maar moet zij het teken van $F$
veranderen: $F = -F$, dus $F = 0$ daar. \omterm{def:b1:arith:divides}{Deelbaarheid}: vat $F$ op
als \omterm{def:b1:poly:def}{veelterm} in de enkele veranderlijke $x_n$ met coëfficiënten in
de overige veranderlijken; zij wordt nul in de $n - 1$ ``waarden''
$x_1, \dots, x_{n-1}$, dus geeft herhaald afsplitsen
(\cref{thm:b1:poly:factor}) dat $F = \prod_{i<n}(x_n - x_i)\cdot G$
met $G$ een \omterm{def:b1:poly:def}{veelterm}. De voorfactor is invariant onder
verwisselingen van twee indices $i, j < n$, dus is $G$ alternerend
in $x_1, \dots, x_{n-1}$, en de inductie voltooit het bewijs:
$\prod_{i<j}(x_j - x_i)$ \omterm{def:b1:arith:divides}{deelt} $F$.

\textbf{24.} $D = \det(x_i^{\,j-1})$ is een \omterm{def:b1:poly:def}{veelterm} in de $x_i$;
twee veranderlijken verwisselen verwisselt twee rijen, dus is $D$
alternerend, en volgens vraag 23 is $D = c\,\prod_{i<j}(x_j - x_i)$
voor een zekere \omterm{def:b1:poly:def}{veelterm} $c$. Totale graden: $D$ heeft graad $\leq 0
+ 1 + \dots + (n-1) = \binom n2$, en het product heeft graad precies
$\binom n2$: $c$ is een constante. Het monoom $x_2\,x_3^2\cdots
x_n^{\,n-1}$ heeft coëfficiënt $1$ in $D$ (het diagonale product) en
$1$ in het product (kies in elke factor de veranderlijke met de
grootste index): $c = 1$, en de formule van Vandermonde valt eruit
zonder enige inductie.

\textbf{25.} (i) Het alterneren --- het axioma ``twee gelijke
kolommen doden de determinant'' --- is de motor: het bracht elke
factor $(x_j - x_i)$, $(a_j - a_i)$, $(b_j - b_i)$ van deze opgave
voort. (ii) De identiteit $\det S = V^2$ vervangt de afzonderlijk
complexe, onbereikbare wortels door hun machtsommen, die reële
\omterm{def:b1:poly:def}{veeltermen} in de coëfficiënten zijn, zodat de onderscheidbaarheid
het teken van een berekenbaar reëel getal wordt. (iii) De
\omterm{ex:b1:det:vandermonde}{Vandermonde-determinant} bestuurt de interpolatie met \omterm{def:b1:poly:def}{veeltermen}, en
de \omterm{pb:b1:det:1}{Cauchy-determinant} bestuurt de partieelbreuken en de bemonsterde
rationale functies (met de \omterm{pb:b1:det:1}{Hilbert-matrix} als beroemdste bijzondere
geval). (iv) Elke alternerende \omterm{def:b1:poly:def}{veelterm} is \omterm{def:b1:arith:divides}{deelbaar} door
$\prod_{i<j}(x_j - x_i)$, en een graadtelling pint zo'n \omterm{def:b1:poly:def}{veelterm}
daarna op een constante na vast --- en dat is waarom dit product
overal blijft opduiken waar een determinant op samenvallingen nul
wordt. De stelling van deel III is de \emph{dubbele alternant van
Cauchy}.
\end{solution}
